\begin{figure}[H]
    \centering
    \hrule
    \vspace{0.3cm}
    \begin{subfigure}[b]{0.48\textwidth}
        \centering
        \small{State $s_1$}
        \begin{tikzpicture}[scale=1.3]
            % Draw the grid and axes
            % \draw[step=0.5cm,very thin,mylightgrey] (0,0) grid (3,3);
            \draw[thin,->] (0,0) -- (3.2,0) node[right] {$q(s_1,\best{a})$};
            \draw[thin,->] (0,0) -- (0,2.2) node[above] {$q(s_1,\worst{a})$};

            % \node at (1.6, 2.2) [above, font=\small, color=black] {State $s_1$};

            % Add labels to the axes
            % \draw (0,0.1) -- (0,-0.1) node[left, font=\small, xshift=-4.5pt, yshift=-6.5pt] {0};
            \draw (0,0.1) -- (0,-0.1) node[below, font=\small] {0};
            \draw (-0.1,0) -- (0.1,0) node[font=\small] {};
            \foreach \x in {1,2,3}
                \draw (\x,0.1) -- (\x,-0.1) node[below, font=\small] {\x};
            \foreach \y in {1,2}
                \draw (0.1,\y) -- (-0.1,\y) node[left, font=\small] {\y};
            
            % Draw the grey square from (1,1) to (2,2)
            \filldraw[fill=mylightgrey, draw=none] (1,1) rectangle (2,2);

            % Draw the line y = x
            \draw[dotted, black] (0,0) -- (2.2,2.2);

            % Draw the blue triangle
            % \filldraw[pattern=dots, pattern color=CambridgeCoreBlue, pattern size=10pt, draw=CambridgeCoreBlue, opacity=1] (1,1) -- (2,2) -- (2,1) -- cycle;
            \filldraw[fill=CambridgeSuperLightBlue, draw=CambridgeCoreBlue, opacity=1] (1,1) -- (2,2) -- (2,1) -- cycle;
            
            % Add the points with associated text
            \filldraw (3,0) circle (1pt) node[above] {$q_{\pol_1}$};
            \filldraw[fill=CambridgeCoreBlue, draw=CambridgeCoreBlue] (1,1) circle (1pt) node[above left, text=CambridgeCoreBlue] {$q_{\pol_5}$};
            \filldraw[fill=CambridgeCoreBlue, draw=CambridgeCoreBlue] (1,1) circle (1pt) node[below right, text=CambridgeCoreBlue] {$q_{\pol_4}$};
            \filldraw[fill=CambridgeCoreBlue, draw=CambridgeCoreBlue] (2,2) circle (1pt) node[right, text=CambridgeCoreBlue] {$q_{\pol_3}$};
            \filldraw (2,0) circle (1pt) node[above] {$q_{\pol_2}$};
            \filldraw (1,0) circle (1pt) node[above right] {$q_{\pol_6}$};
            \filldraw (1,0) circle (1pt) node[above left] {$q_{\worst{\pol}}$};

            \node at (1.5, 1.5) [font=\small, color=black] {$\qset_0$};
        \end{tikzpicture}
        \caption{$q_k(s_1, \best{a}) > q_k(s_1, \worst{a})$}
        % \caption{$\best{a} = \pol_3(s_1) = \pol_4(s_1) = \pol_5(s_1)$}
        \label{fig: explain stay in box for loop action best}
    \end{subfigure}
    \hfill
    \begin{subfigure}[b]{0.48\textwidth}
        \centering
        \small{State $s_1$}
        \begin{tikzpicture}[scale=1.3]
            % Draw the grid and axes
            % \draw[step=0.5cm,very thin,mylightgrey] (0,0) grid (3,3);
            \draw[thin,->] (0,0) -- (3.2,0) node[right] {$q(s_1,\best{a})$};
            \draw[thin,->] (0,0) -- (0,2.2) node[above] {$q(s_1,\worst{a})$};

            % \node at (1.6, 2.2) [above, font=\small, color=black] {State $s_1$};

            % Add labels to the axes
            % \draw (0,0.1) -- (0,-0.1) node[left, font=\small, xshift=-4.5pt, yshift=-6.5pt] {0};
            \draw (0,0.1) -- (0,-0.1) node[below, font=\small] {0};
            \draw (-0.1,0) -- (0.1,0) node[font=\small] {};
            \foreach \x in {1,2,3}
                \draw (\x,0.1) -- (\x,-0.1) node[below, font=\small] {\x};
            \foreach \y in {1,2}
                \draw (0.1,\y) -- (-0.1,\y) node[left, font=\small] {\y};
            
            % Draw the grey square from (1,1) to (2,2)
            \filldraw[fill=mylightgrey, draw=none] (1,1) rectangle (2,2);

            % Draw the line y = x
            \draw[dotted, black] (0,0) -- (2.2,2.2);

            % Draw the blue triangle
            \filldraw[fill=CambridgeSuperLightBlue, draw=CambridgeCoreBlue, opacity=1] (1,1) -- (2,2) -- (1,2) -- cycle;
            
            % Add the points with associated text
            \filldraw[fill=CambridgeCoreBlue, draw=CambridgeCoreBlue] (3,0) circle (1pt) node[above, text=CambridgeCoreBlue] {$q_{\pol_1}$};
            \filldraw (1,1) circle (1pt) node[above left] {$q_{\pol_5}$};
            \filldraw (1,1) circle (1pt) node[below right] {$q_{\pol_4}$};
            \filldraw (2,2) circle (1pt) node[right] {$q_{\pol_3}$};
            \filldraw[fill=CambridgeCoreBlue, draw=CambridgeCoreBlue] (2,0) circle (1pt) node[above, text=CambridgeCoreBlue] {$q_{\pol_2}$};
            \filldraw (1,0) circle (1pt) node[above left] {$q_{\worst{\pol}}$};
            \filldraw[fill=CambridgeCoreBlue, draw=CambridgeCoreBlue] (1,0) circle (1pt) node[above right, text=CambridgeCoreBlue] {$q_{\pol_6}$};
            

            \node at (1.5, 1.5) [font=\small, color=black] {$\qset_0$};        
        \end{tikzpicture}
        \caption{$q_k(s_1, \best{a}) < q_k(s_1, \worst{a})$}
        \label{fig: explain stay in box for loop action worst}
    \end{subfigure}%  
    \vspace{-0.2cm}
    \caption{
    % Whether action $\best{a}$ or $\worst{a}$ is greedy in states $s_1$, $s_2$ or $s_3$,
    {\sffamily\textbf{(a)}} When $\best{a}$ is greedy in $s_1$, solutions move towards $q_{\pol_3}$, $q_{\pol_4}$ or $q_{\pol_5}$.
    {\sffamily\textbf{(b)}} When $\worst{a}$ is greedy in $s_1$, solutions move towards $q_{\pol_1}$, $q_{\pol_2}$ or $q_{\pol_6}$.
    Either way, 
    solutions starting in $\qset_0$ remain in $\qset_0$ when learning rates are small.
    % and it only visits the policies $\pol_1$ through $\pol_6$.
    This conclusion also holds in states $s_2$ and $s_3$.
    }
    \label{fig: explain stay in box for loop}
\end{figure}